% Avsnitt 5.7, ur lectures/L05/appendix/b_exercises.md.

\section{Övningar}
\label{c5:sec:exercises}

\begin{aside}
\textbf{Så kontrollerar du ditt arbete.} Varje övning nedan som ber dig skriva en VHDL-modul har en
självkontrollerande testbänk under \file{lectures/L05/exercises/}. Skriv din modul i dess katalog
\file{exercises/<module>/}, med det entitetsnamn och den \textbf{portordning} övningen anger, och
kör den sedan med GHDL \textemdash{} se \secref{c2:sec:testbenches} för de tre kommandona.
\end{aside}

% ------------------------------------------------------------------------------------------
\exerciseset{\code{variable} kontra \code{signal}}

\exercise{Spåra ackumulatorbuggen}{Resonemang}
\label{c5:ex:trace}
Den trasiga \code{parity_gen}-arkitekturen från \secref{c5:sec:variable} återges nedan:

\begin{vhdlcode}
architecture broken of parity_gen is
signal parity_s: std_logic;
begin
    parity <= parity_s;

    process(bits) is
    begin
        parity_s <= '0';
        for i in bits'range loop
            parity_s <= parity_s xor bits(i);
        end loop;
    end process;
end architecture;
\end{vhdlcode}

Anta att \code{bits} har intervallet \code{7 downto 0}, att \code{bits = "10110110"}, och att
\code{parity_s} håller \code{'0'} innan processen börjar.

\xpart{a} Spåra processen för hand. Notera, för den inledande tilldelningen och för varje
loopiteration: det aktuella indexet \code{i}, värdet på \code{bits(i)}, värdet som läses från
\code{parity_s}, värdet som schemaläggs för \code{parity_s}, och om den schemalagda tilldelningen
överlever fram till att processen suspenderar.

Kom ihåg att en signaltilldelning inte uppdaterar signalen omedelbart, och att tilldelningar till
samma signal under ett och samma varv ersätter tidigare schemalagda tilldelningar för samma
simuleringstidpunkt.

\xpart{b} Bestäm det slutliga värde som tilldelas \code{parity_s} efter att processen suspenderat.
Beräkna sedan den korrekta XOR-pariteten hos \code{"10110110"} för hand, jämför den med resultatet
från den trasiga arkitekturen, och förklara varför upprepade tilldelningar till signalen inte
skapar en ackumulator.

\xpart{c} Spåra den variabelbaserade arkitekturen \code{behaviour} från \file{parity_gen.vhd} med
samma ingång. Notera, för varje loopiteration: värdet på \code{bits(i)}, värdet på \code{acc} före
XOR-operationen, och värdet på \code{acc} omedelbart efter den. Bekräfta att \code{acc} behåller
varje iterations resultat så att nästa iteration kan använda det, och att det slutliga värde som
tilldelas \code{parity} är den korrekta XOR-pariteten.

\exercise{\code{min_of_two}}{VHDL}
\label{c5:ex:minoftwo}
Skriv en entitet med namnet \code{min_of_two}.

\begin{booktable}{@{}p{3.6em}p{3.4em}p{11.4em}L@{}}
  \thead{Port} & \thead{Riktn.} & \thead{Typ} & \thead{Betydelse} \\
  \midrule
  \code{a} & in & \code{std_logic_vector(3 downto 0)} & Första värdet. \\
  \code{b} & in & \code{std_logic_vector(3 downto 0)} & Andra värdet. \\
  \code{m} & out & \code{std_logic_vector(3 downto 0)} & Det minsta av de två. \\
\end{booktable}

Behandla \code{a} och \code{b} som teckenlösa fyrabitarsvärden. Implementera kretsen med en process:
\begin{itemize}
  \item Ta med \code{a} och \code{b} i känslighetslistan.
  \item Deklarera en processlokal variabel som håller det minsta värdet.
  \item Jämför \code{a} och \code{b} som teckenlösa tal, med faciliteterna i
    \code{ieee.numeric_std} enligt \secref{c2:sec:vectors}.
  \item Tilldela variabeln exakt en gång per varv, och tilldela den till \code{m} efter
    jämförelsen.
\end{itemize}

Den här övningen kräver ingen loop. Poängen med den är att använda en variabel som tillfällig
processlokal lagring för något \emph{annat} än en ackumulator: \secref{c5:sec:variable}:s
genomarbetade exempel bygger upp ett värde över iterationerna i en loop, vilket är det fall en
variabel oftast introduceras för, och det är lätt att gå därifrån och tro att det är det enda de
duger till. Här håller variabeln helt enkelt ett mellanresultat under de få rader som skiljer
beräkningen från användningen, vilket är den vanligare av de två användningarna i verklig kod.

\modulefig[0.62]{lectures/L05/appendix/images/min_of_two.png}

\begin{selfcheck}
\textbf{Självkontroll.} Döp din entitet till \code{min_of_two}, med portarna \code{a}, \code{b}
(in) och \code{m} (out), alla \code{std_logic_vector(3 downto 0)}, deklarerade i den ordningen;
dess testbänk finns i \file{exercises/min_of_two/}.
\end{selfcheck}

\exercise{\code{ones_count8}}{VHDL}
\label{c5:ex:ones}
Skriv en entitet med namnet \code{ones_count8} som räknar hur många bitar i en 8-bitars ingång som
är ettställda.

Övning~\ref{c5:ex:trace} lät dig spåra ackumulatorbuggen på papper, och Övning~\ref{c5:ex:minoftwo}
använde en variabel som inte var någon ackumulator alls. Den här är det fall
\secref{c5:sec:variable} faktiskt handlar om: ett värde som byggs upp över iterationerna i en loop.
Det är den enda övningen i boken där du skriver en sådan.

\begin{booktable}{@{}p{4em}p{3.4em}p{11.4em}L@{}}
  \thead{Port} & \thead{Riktn.} & \thead{Typ} & \thead{Betydelse} \\
  \midrule
  \code{bits} & in & \code{std_logic_vector(7 downto 0)} & Värdet att räkna över. \\
  \code{ones} & out & \code{natural range 0 to 8} & Antalet bitar i \code{bits} som är \code{'1'}. \\
\end{booktable}

Implementera den med en enda kombinatorisk process:
\begin{itemize}
  \item Sätt \code{bits} i känslighetslistan.
  \item Deklarera en processlokal \code{variable} för den löpande räkningen, och sätt den till
    \code{0} överst i varje varv. En variabel behåller sitt värde mellan varv, så en räkning som
    aldrig nollställs fortsätter växa.
  \item Loopa över \code{bits'range} och lägg till \code{1} för varje ettställd bit.
  \item Tilldela variabeln till \code{ones} en gång, efter loopen.
\end{itemize}

Svara sedan, med en mening var:
\begin{itemize}
  \item Skriv samma sak med en \code{signal} som ackumulator i stället, kör testbänken, och notera
    vad den rapporterar för ingången \code{00000010}. Förklara talet du får med
    schemaläggningsregeln, inte genom att gissa.
  \item Loopen körs åtta gånger. Hur många adderare lägger den här designen på FPGA:n, och hur
    många klockcykler tar det för den att producera ett svar? Om de två svaren överraskar dig, läs
    om \secref{c5:sec:variable}.
  \item \code{ones} är en \code{natural range 0 to 8} snarare än ett
    \code{std_logic_vector(3 downto 0)}. Båda rymmer svaret. Vad köper intervallet dig, och vad
    skulle GHDL göra om din räkning någonsin lämnade det?
\end{itemize}

\modulefig[0.66]{lectures/L05/appendix/images/ones_count8.png}

\begin{selfcheck}
\textbf{Självkontroll.} Döp din entitet till \code{ones_count8}, med ingången \code{bits}
(\code{std_logic_vector(7 downto 0)}) och utgången \code{ones} (\code{natural range 0 to 8}),
deklarerade i den ordningen; dess testbänk finns i \file{exercises/ones_count8/} och sveper igenom
alla 256 ingångsvärden.
\end{selfcheck}

% ------------------------------------------------------------------------------------------
\exerciseset{En modul värd att ha för sig själv}

\exercise{\code{sync}: synkroniseraren som egen modul}{VHDL}
\label{c5:ex:sync}
Skriv dubbelvippsynkroniseraren som en egen modul, med namnet \code{sync}.

Du har redan byggt den här kretsen. Det är de två första vipporna i \chapref{c4:ch}:s
\code{button_sync}, skrivna inline där eftersom det vid den punkten var \emph{idén} som spelade
roll: en asynkron ingång behöver två vippor innan något annat får titta på den. Här bryter du ut de
två till en modul, och övningen är inte logiken, som du redan kan, utan allt runt omkring.

Två saker gör det värt att ta en gång till. Den första är att en naken synkroniserare är vad de
flesta asynkrona ingångar faktiskt behöver: \code{button_sync} lägger till en flankdetektor
eftersom en knapptryckning är en händelse, men en seriell ledning, en statusflagga från en annan
klockdomän, eller en strömbrytare vars nivå är allt du någonsin läser, behöver de två vipporna och
inget mer. \Chapref{c8:ch}:s seriemottagare instansierar precis det här på sin \code{rx}-pinne. Den
andra är att den tar \textbf{två generics}, och en av dem är inte ett tal.

\begin{booktable}{@{}p{4.4em}p{9.8em}p{4em}L@{}}
  \thead{Generic} & \thead{Typ} & \thead{Förval} & \thead{Betydelse} \\
  \midrule
  \code{SIZE} & \code{natural range 1 to 15} & \code{1} & Hur många oberoende signaler som ska
    synkroniseras, och därmed bredden på båda vektorportarna. \\
  \code{PRESET} & \code{std_logic} & \code{'0'} & Värdet som båda stegen antar vid reset. \\
\end{booktable}

och de här portarna:

\begin{booktable}{@{}p{5.6em}p{3.4em}p{11em}L@{}}
  \thead{Port} & \thead{Riktn.} & \thead{Typ} & \thead{Betydelse} \\
  \midrule
  \code{clock} & in & \code{std_logic} & 50 MHz systemklocka. \\
  \code{reset_s2_n} & in & \code{std_logic} & Aktiv låg, \textbf{redan synkroniserad} reset. En
    synkroniserare synkroniserar inte sin egen reset; det gör \chapref{c4:ch}:s
    \code{reset_sync}. \\
  \code{async_in} & in & \code{std_logic_vector(SIZE - 1 downto 0)} & Den asynkrona ingången eller
    ingångarna, rakt utifrån klockdomänen. \\
  \code{sync_out} & out & \code{std_logic_vector(SIZE - 1 downto 0)} & Samma signaler, säkra att
    använda, två stigande flanker senare. Varje bit synkroniseras oberoende. \\
\end{booktable}

\xpart{a} Implementera den:
\begin{itemize}
  \item Två interna signaler av typen \code{std_logic_vector(SIZE - 1 downto 0)}, en per steg.
  \item En process, känslig för \code{clock} och \code{reset_s2_n}.
  \item Vid reset, driv båda stegen till \code{PRESET}, utanför den klockade grenen.
  \item Vid en stigande klockflank, skifta \code{async_in} genom de två stegen.
  \item Driv \code{sync_out} från det andra steget med en konkurrent tilldelning.
\end{itemize}

\xpart{b} \code{PRESET} är en generic av typen \code{std_logic}, den första i den här boken som
inte är ett tal, och den finns eftersom det säkra reset-värdet beror på vad signalen \emph{betyder}.
Bara den instansierande designen vet det. Svara:
\begin{itemize}
  \item \code{button_sync} synkroniserar aktivt låga knappar, så den skulle skicka in \code{'1'}.
    Vad skulle \code{'0'} påstå om knapparna under de två cyklerna efter varje reset?
  \item En seriell ledning vilar hög, så \chapref{c8:ch}:s mottagare skickar också in \code{'1'}.
    Hur skulle \code{'0'} se ut för en mottagare som håller utkik efter en startbit?
  \item Varför är \emph{standardvärdet} ändå \code{'0'} snarare än \code{'1'}, givet båda dessa,
    och vad säger det om vem som ansvarar för att få det rätt?
\end{itemize}

\xpart{c} Det här är också en fråga om vad en generic kostar. \code{SIZE} är
\code{natural range 1 to 15} och \code{PRESET} är \code{std_logic}, och båda är fastlagda innan
syntesen körs. Använd \secref{c5:sec:fabric}:s bild av vad FPGA:n faktiskt består av och säg hur
många vippor en instans med \code{SIZE = 4} använder, hur många en instans med \code{SIZE = 1}
använder, och var i den byggda kretsen \code{PRESET} hamnar. Ingen av de två dyker upp som logik
någonstans; säg vad som hände med dem i stället.

\xpart{d} Verifiera den med dess testbänk. Den instansierar din modul två gånger, en gång en bit
bred med \code{PRESET = '1'} och en gång tre bitar bred med \code{PRESET = '0'}, så en design som
hårdkodar någon av de två generics fallerar. Den kontrollerar att \code{sync_out} följer
\code{async_in} efter \textbf{exakt} två stigande flanker snarare än efter en eller så småningom,
att varje bit i en vektor synkroniseras för sig, och att reset laddar \code{PRESET} utan att någon
klockflank är inblandad.

\note[Tips]Skriv arkitekturen för \code{SIZE = 1} i huvudet först, och kontrollera sedan att varje
rad redan är korrekt för en vektor. Aggregatet \code{(others => PRESET)} och en vanlig
vektortilldelning arbetar båda elementvis, så generaliseringen bör inte kosta någon extra kod.

\modulefig[0.76]{lectures/L05/appendix/images/sync.png}

\begin{selfcheck}
\textbf{Självkontroll.} Döp din entitet till \code{sync}, med generics \code{SIZE}
(\code{natural range 1 to 15}, standardvärde \code{1}) och \code{PRESET} (\code{std_logic},
standardvärde \code{'0'}), ingångarna \code{clock}, \code{reset_s2_n}, \code{async_in}
(\code{std_logic_vector(SIZE - 1 downto 0)}) och utgången \code{sync_out}
(\code{std_logic_vector(SIZE - 1 downto 0)}), deklarerade i den ordningen; dess testbänk finns i
\file{exercises/sync/}.

\textbf{Spara den här filen.} \Chapref{c8:ch}:s capstone med seriemottagaren instansierar den på
\code{rx}-pinnen, och ber dig kopiera in den.
\end{selfcheck}

% ------------------------------------------------------------------------------------------
\exerciseset{Från VHDL till hårdvara}

Quartus-flödet demonstreras på ett riktigt DE0-CV-kort under passen \textemdash{} först med
bromsassistenten redan i \chapref{c1:ch} \textemdash{} och finns nedskrivet i
bilaga~\ref{app:quartus} som referens. Du behöver varken Quartus eller ett kort själv; övningarna
nedan handlar om att förstå vad det flödet gör med den VHDL du skriver.

\exercise{Vad verktygskedjan gör}{Resonemang}
\label{c5:ex:quartus}
Besvara vart och ett av följande, med hänvisning till bilaga~\ref{app:quartus} och till det du såg
demonstreras. En eller två meningar var.

\xpart{a} Ett Quartus-projekt behöver en vald \textbf{toppnivåentitet} innan det kompilerar. Vad
gör verktyget annorlunda med toppnivåentiteten än med någon annan entitet i projektet?
\code{min_of_two} från Övning~\ref{c5:ex:minoftwo} och dess testbänk är båda giltig VHDL. Varför
kan bara den ena någonsin vara toppnivåentitet i ett FPGA-projekt?

\xpart{b} Pinntilldelning kopplar ett portnamn till en fysisk pinne. Varför måste pinnarna
tilldelas \emph{före} kompileringen snarare än efteråt? Vad skulle \code{m} fysiskt vara kopplad
till om du kompilerade \code{min_of_two} utan att ha tilldelat några pinnar alls?

\xpart{c} Syntesen gör om din arkitektur till grindar på chippet. \Chapref{c2:ch}:s övning
\code{combo_logic} lät dig skriva $x = ab + c'd$ två gånger: en gång som ett enda uttryck, och en
gång med $c'd$ namngivet av en intern signal. Vad bör syntesen producera för vardera, och varför är
ditt svar detsamma för båda? Vilken av de två skulle du hellre läsa om sex månader?

\xpart{d} Kompileringen ger varningar såväl som fel. Varför är ”det kompilerade utan fel” ett mycket
svagare påstående om en FPGA-design än ”det kompilerade utan fel” är om ett C-program? Nämn en sorts
varning du aldrig skulle ignorera.

\exercise{Fyra förutsägelser om pariteten}{Simulering}
\label{c5:ex:parity}
\code{parity_gen} från \secref{c5:sec:variable} är kapitlets genomarbetade exempel, och till det hör
en referenstestbänk.

\xpart{a} Förutsäg, innan du kör något, \code{parity} för hand för var och en av de här ingångarna:
\code{"00000000"}, \code{"10110110"}, \code{"11111111"} och \code{"10000000"}.

\xpart{b} Kör referenstestbänken mot det genomarbetade exemplet och bekräfta dina fyra
förutsägelser:

\begin{shell}
cd lectures/L05/parity_gen
ghdl -a --std=93 parity_gen.vhd parity_gen_tb.vhd
ghdl -e --std=93 parity_gen_tb
ghdl -r --std=93 parity_gen_tb --assert-level=error --stop-time=10ms
\end{shell}

\xpart{c} Ta vilken som helst av de ingångarna och vänd exakt en bit, och lämna de andra sju
oförändrade. Vad händer med \code{parity}? Förklara varför en ändring av exakt en ingångsbit
\emph{alltid} måste invertera XOR-paritetsresultatet, vilken bit du än väljer och vilket
startvärdet än är. Det är den egenskapen som gör en paritetsbit förmögen att upptäcka ett
enbitsfel vid överföring. Vilken sorts fel missar den, och varför?

\exercise{En avkodare som kompilerar rent och ändå är fel}{Resonemang}
\label{c5:ex:latch}
En kollega skickar dig den här kombinatoriska avkodaren för granskning. Den kompilerar, och Quartus
rapporterar inga fel.

\begin{vhdlcode}
architecture behaviour of level_decode is
begin
    process(sel) is
    begin
        case sel is
            when "00"   => leds <= "0001";
            when "01"   => leds <= "0010";
            when "10"   => leds <= "0100";
            when others => null;
        end case;
    end process;
end architecture;
\end{vhdlcode}

\xpart{a} Quartus ger \code{Warning: Inferred latch(es) for signal "leds"}. Förklara med hjälp av
\secref{c5:sec:latch} vad verktyget byggde och varför. Vad bad VHDL-koden om som ett kombinatoriskt
nät inte kan göra?

\xpart{b} Författarens försvar är att \code{sel} är två bitar, så \code{"11"} är det enda värde
\code{when others} kan matcha, och designen genererar det aldrig. Ge två skilda skäl till att det
argumentet inte håller: ett om vad \code{sel} faktiskt kan hålla (\secref{c1:sec:vhdl}), och ett
som skulle gälla även om \code{sel} verkligen bara kunde vara \code{"00"}, \code{"01"} eller
\code{"10"}.

\xpart{c} Åtgärda den, på två olika sätt: genom att ändra enbart grenen \code{when others}, och
genom att lägga till en enda rad \emph{före} \code{case}-satsen och lämna varje gren orörd. Vilken
av de två skulle du hellre underhålla, och varför? Den andra skalar till ett \code{case} med tjugo
grenar; den första gör det inte.

\xpart{d} Det här är en varning, inte ett fel, och designen kommer ofta att verka fungera på
kortet. Varför gör det saken farligare än ett fel snarare än mindre farlig?
